\recipe{Pietro’s Pasta alla Siciliana}
\label{recipe:pasta_alla_siciliana}
\makes{x portions}
\worktime{\SI{30} \minute}
\totaltime{\SI{40} \minute}
\source{Pietro}

A bold Southern Italian pasta with a rich anchovy-tomato sauce, sweet raisins, crunchy toasted breadcrumbs and pine nuts


\begin{ingredients} % INGREDIENTI

  \item 500 g thick spaghetti 
  \item 200 g concentrated tomato paste 
  \item 10 anchovies in olive oil 
  \item 40 g pine nuts
  \item 50 g raisins (uvetta) 
  \item 100 g onion, finely chopped 
  \item 50 g shallot (scalogno), finely chopped 
  \item 1 garlic clove 
  \item 100 g breadcrumbs 
  \item 59 1/8 ml olive oil 
  \item 1 pinch black pepper
\end{ingredients}


\begin{method} % Passi
  	 Heat a generous amount of olive oil in a pan over medium heat. 
Add the 1 garlic clove clove and fry until golden, then remove and
discard it. 

Add the 100 g onion, finely chopped and 50 g shallot (scalogno), finely
chopped and cook until golden and soft.

Add the 10 anchovies in olive oil to the pan and stir until they completely
dissolve into the oil and onions. 

Stir in the 50 g raisins (uvetta) and let them cook for a few minutes so
they plump up and absorb the flavours. 

Add the 200 g concentrated tomato paste and a splash of water. Stir well
and simmer on low heat until the sauce comes together. Taste and adjust
salt if needed (anchovies are already salty). (10m) 

In a separate dry or lightly oiled pan, toast the 100 g breadcrumbs with a
drizzle of olive oil over medium heat, stirring frequently, until golden
brown and crispy. 

Set aside. (5m) 6. In the same pan, toast the 40 g pine nuts with a little
olive oil until golden and crunchy. Add a pinch of 1 pinch black pepper
towards the end. The crunchier the better. Set aside. (4m) 

Bring a large pot of salted water to a boil and cook the 500 g thick
spaghetti until al dente according to the package instructions. Reserve a
cup of pasta water before draining. (10m) 

Toss the drained pasta with the sauce, adding a splash of pasta water if
needed to loosen it. Plate and top generously with the toasted
breadcrumbs and pine nuts. Serve immediately
  
\end{method}


\begin{figure}[h]
  	\centering
  	
  	\begin{tikzpicture}
  		% Save the image in a box to get its dimensions
  		\sbox0{\includegraphics[width = 0.6\linewidth]{images/MissingImage.png}}%
  
  		% Draw a clipping rectangle with rounded corners
  		\path[clip, rounded corners=0.5cm] (0,0) rectangle (\wd0,\ht0);
  
  		% Place the image inside the clipped area
  		\node[inner sep=0pt, outer sep=0pt, anchor=south west] at (0,0) {\usebox0};
	\end{tikzpicture}

  \caption*{caption for image with date and information, 1st January 2026}
\end{figure}

\begin{notes}
Tips: Use thick spaghetti (spaghettoni) for the
best result — they hold up to the bold
sauce. The toasted breadcrumbs
(pangrattato) are the Sicilian substitute for
grated cheese, so don't skip them!
Anchovies are salty, so taste before adding
any extra salt.

\end{notes}

